\documentclass[12pt,a4paper]{article}
\usepackage[a4paper,landscape,top=2cm,bottom=2cm,left=2cm,right=2cm]{geometry}
\usepackage{fontspec}
\setmainfont{Thonburi}
\XeTeXlinebreaklocale "th"
\XeTeXlinebreakskip=0pt plus 1pt
\usepackage{booktabs,tabularx,array}
\setlength{\parindent}{0pt}
\setlength{\emergencystretch}{3em}
\begin{document}
\begin{center}
\small
\textbf{ตารางที่ 3.13 สรุปผลเชิงตัวเลขและการคัดเลือกแบบจำลองในแต่ละสายการทดลอง}

\vspace{0.5em}
\begin{tabularx}{\linewidth}{p{3.2cm}p{4cm}p{2cm}XX}
\toprule
สายการทดลอง & แบบจำลอง/ชุดคุณลักษณะ & จำนวนคุณลักษณะ & ผล Validation หรือเกณฑ์คัดเลือก & ผลการทดสอบ (Test) \\
\midrule
1. Classical Leakage-safe & Gradient Boosting & 22 & Average Precision = 0.9733\newline F1 = 0.9368 & Average Precision = 0.9703\newline F1 = 0.9337 \\
\addlinespace
2. Multimodal Suite (Final Model) & Gradient Boosting (\texttt{image\_context}) & 41 & Composite Score = 0.9613\newline Average Precision = 0.9776\newline F1 = 0.9339 & Average Precision = 0.9789\newline F1 = 0.9340 \\
\addlinespace
3. Neural Reference & IdentityMLP & 22 & Average Precision $\approx$ 1.0000\textsuperscript{1} & Average Precision = 0.7344\newline F1 = 0.8296 \\
\bottomrule
\end{tabularx}
\end{center}

\footnotesize
\textsuperscript{1} ค่า Average Precision $\approx$ 1.0000 เป็นค่าที่ปรากฏในรายงานเดิมของสาย Neural Reference แต่ไม่ได้ระบุชุดข้อมูลไว้อย่างชัดเจน จึงไม่ควรนำมาเปรียบเทียบกับผล Test โดยตรง ส่วนผล Test ที่ตรวจสอบได้คือ Average Precision = 0.7344 และ F1 = 0.8296

\textbf{หมายเหตุ:} \texttt{image\_context} เป็นชื่อชุดคุณลักษณะของสาย Multimodal ส่วนแบบจำลองที่ใช้คือ Gradient Boosting
\end{document}
